\section{Related Work}
\subsection{Interactive Long Video Generation}
In following parts, we will review related work from four aspects:

\paragraph{Video Generation Model.} Video generation models can generate video sequences $\mathbf{x} = \{ x^0, x^1, \dots, x^t\}$, where $x^i$ indicates the $i$-th frame. The current mainstream model architecture is based on diffusion models~\cite{ddpm, sde, score,flow,rectified}, which excel in generating high-quality content and have been widely adopted in video generation~\cite{sora,cogvideox,kling,veo2,runway,vidu,wan,hunyuanvideo}. Other alternative architectures include next-token prediction~\cite{videogpt,videopoet,emu3} and various hybrid approaches~\cite{cdf,nova,mar}.

\paragraph{Controllable Video Generation.} This task can be formulated as $p(\mathbf{x}|c)$, where $c$ represents different types of control signals. The most representative control signals include: camera motion control~\cite{motionctrl,cameractrl,3dtrajmaster,recammaster, syncammaster}, and agent action control in games or simulators~\cite{gamengine,genie2,gamefactory,matrix,oasis}. These control signals greatly enhance user interactive experience, enabling free exploration in the created virtual worlds.

\paragraph{Streaming Video Generation.} Streaming video generation can condition on previously generated frames to continuously generate new video frames, which can be expressed as $p(x^0,x^1,...,x^n)=\prod_{i=0}^np(x^i|x^0,x^1,...,x^{i-1})$, where $x^i$ indicates the $i$-th frame. Representative approaches include Diffusion-based methods~\cite{cdf,dfot,gamefactory,far,framepack} and GPT-like next token prediction methods~\cite{videopoet,emu3,wham}. Diffusion-based methods generally achieve higher visual quality and faster sampling speed, thus we focus on diffusion models for long video generation in this work. Although these SOTA methods generally fail to generate long videos with scene-consistent memory, instead only producing long videos with short-term continuity.


\paragraph{Memory Capability for Video Generation.} Many related works' demos~\cite{gamengine,oasis,dfot,wham} have shown that current long video generation methods generally lack memory capability: while maintaining frame-to-frame continuity, the scenes continuously change. One potential approach~\cite{wonderjourney,viewcrafter,see3d,gen3c} is to leverage 3D reconstruction to build explicit 3D representations from generated videos, then render initial frames from these 3D representations as conditions for new video generation. However, this method is limited by the accuracy and speed of 3D reconstruction, particularly in continuously expanding large scenes where accumulated 3D reconstruction errors become intolerable. Moreover, these works focus on 3D generation and merely borrow priors from video generation models, which differs from our scope. WorldMem~\cite{worldmem} attempts to implement memory by injecting historical frames through cross attention, and has been validated on video lengths of around $10$ seconds in Minecraft scenarios.


\subsection{Context Learning for Video Generation}
Recently, some works~\cite{lct,far,framepack} have begun to explore the role of long-context in video generation. LCT~\cite{lct} performs long-context tuning on pre-trained single-shot video diffusion models to achieve consistency in multi-shot video generation. FAR~\cite{far} proposes Long-Term and Short-Term context windows to condition video generation models for long video generation. FramePack~\cite{framepack} introduces a hierarchical method to compress context frames into a fixed number of frames as conditioning for video generation models to achieve long video generation. However, their compression method loses too much information from temporally distant frames. In this work, we further highlight the significance of context, emphasizing that all history context serves as memory for scene-consistent long video generation.



